% Methods fragment for the peripersonal-space audio--tactile task.
% Suggested packages in the parent manuscript:
% \usepackage{booktabs}
% \usepackage{xurl} % or \usepackage{url}; xurl gives better URL line breaks
% \usepackage{natbib} % or another citation package compatible with \citep
% Bibliography entries are provided in methods_references.bib.

\section{Methods}

\subsection{Audio--tactile peripersonal-space task}

The experiment used an audio--tactile reaction-time paradigm designed to probe the
spatial boundary of peripersonal space (PPS), following the logic of dynamic
PPS tasks in which tactile processing is facilitated by task-irrelevant sounds
approaching the body \citep{Canzoneri2012DynamicSounds,Serino2015BodyPart}.
On each trial, participants heard a spoken breathing instruction followed by a
binaural looming sound. A brief vibrotactile stimulus was delivered on a subset
of trials, and participants were instructed to respond as quickly as possible
when they detected the tactile stimulus while ignoring the auditory motion. The
critical experimental manipulation was the temporal onset of the tactile pulse
relative to the motion of the looming sound, which was used as a proxy for the
instantaneous auditory-source distance from the stimulated body part. Faster
responses at near than far source positions were interpreted as the behavioural
signature of audio--tactile integration within PPS.

\subsection{Experimental software and reproducibility}

The task was implemented with the Peripersonal Space Toolkit
\citep{Fejer2026PPSToolkit}, which was developed in the course of designing
the present experiment. The toolkit served three functions. First, it generated
and rendered the auditory and tactile stimulus sequences. Second, it controlled
the experimental playback and response logging. Third, it provided a
replication-oriented graphical interface for specifying PPS paradigms in terms
of their experimentally relevant degrees of freedom, such as auditory
trajectory, source spectrum, spatial reference frame, HRIR selection, tactile
site, SOA structure, baseline and catch-trial policy, respiratory or task
phase, block composition, and randomisation constraints. The software therefore
acted both as the implementation of the present task and as a tool for
replicating or modifying related audio--tactile PPS paradigms. Full software
documentation, installation instructions, dependency lists, licenses, and
implementation details are provided in the public repository; the manuscript
describes only the operations that determined the experimental stimuli,
timing, and analysis.

The present protocol specified a repeated two-row structure in which an inhale
or exhale instruction was paired with one of four frontal looming auditory
sources. Each trial lasted 8.000~s: a 4.000~s spoken instruction segment was
followed by a 4.000~s looming-sound segment. The session comprised six blocks,
each containing 46 trials, for 276 trials per participant. Across the full
schedule, each participant received 240 audio--tactile trials and 36 catch
trials. Audio--tactile trials were balanced across inhale and exhale phases,
across five tactile stimulus-onset asynchronies (SOAs), and across four
auditory-source spectra. Catch trials contained the identical auditory sequence
but no tactile pulse. Baseline tactile-only trials were not included in this
protocol, although the software supports such trials for other PPS variants.

The manuscript-facing software attribution was restricted to the components
that performed the methodological heavy lifting (Table~\ref{tab:software-attribution}).

\begin{table}[t]
\centering
\caption{Core software and scientific resources used for stimulus rendering, timing, and analysis.}
\label{tab:software-attribution}
\begin{tabular}{p{0.25\textwidth}p{0.48\textwidth}p{0.19\textwidth}}
\toprule
Component & Methodological role & Citation \\
\midrule
Peripersonal Space Toolkit & Project-developed stimulus generator, experiment player, replication interface, schedule generator, event logger, and analysis-output generator & \citep{Fejer2026PPSToolkit} \\
3D Tune-In Audio Toolkit & Binaural rendering framework for the moving auditory source & \citep{CuevasRodriguez2019ThreeDTI} \\
SOFA/AES69, \texttt{sofar}, and FABIAN & Standardized HRIR format, HRIR loading/inspection, and the non-individualized FABIAN HRIR used for rendering & \citep{Majdak2022SOFA,Sofar2026,Brinkmann2017FABIAN} \\
NumPy and SciPy & Procedural signal generation, numerical signal operations, and reaction-time function fitting & \citep{Harris2020NumPy,Virtanen2020SciPy} \\
\texttt{soundfile}/libsndfile & Audio read/write, stimulus inspection, and multichannel sequence assembly & \citep{SoundFile2026,Libsndfile2026} \\
\texttt{sounddevice}/PortAudio & Multichannel audio playback and device-level audio timing interface & \citep{SoundDevice2026,Bencina2001PortAudio} \\
Lab Streaming Layer & Optional synchronized marker stream for multimodal timing records & \citep{Kothe2025LSL} \\
\bottomrule
\end{tabular}
\end{table}

\subsection{Auditory source generation}

The auditory stimulus was a 4.000~s binaural looming sound sampled at
44.1~kHz. Four auditory source spectra were used: pink, blue, white, and
brown noise. These sounds were first generated as monophonic source signals and
were then rendered as frontal moving binaural sources, producing separate left-
and right-ear signals for headphone presentation.

The virtual auditory source moved along the listener's frontal midline. In the
application coordinate convention, the listener was located at the origin, the
positive \(y\)-axis pointed forward, the positive \(x\)-axis pointed to the
listener's right, and the positive \(z\)-axis pointed upward. The source was
specified at azimuth \(0^\circ\) and elevation \(0^\circ\). Each 4.000~s
looming segment contained a 0.500~s far-position hold, a 3.000~s linear
approach, and a 0.500~s near-position hold. During the approach phase, the
source translated from 110~cm to 10~cm from the listener, corresponding to a
1.000~m path traversed at 33.3~cm/s. The source trajectory therefore
implemented a smooth approach from extrapersonal toward peripersonal space
while preserving a constant frontal direction.

The task schedule used five tactile SOAs, defined as the
delay between the onset of the 4.000~s looming-sound segment and the onset of
the tactile cue: 300, 800, 1500, 2200, and 2700~ms. The protocol labels these
SOA levels with nominal source-distance factors of 100.0, 83.3, 60.0, 36.7,
and 20.0~cm, respectively. The initial 0.500~s far-position hold was included
to avoid an abrupt spatial transition at auditory-stimulus onset and to make
the start of each auditory segment perceptually stable before the main approach
phase. These SOA and distance labels were retained throughout stimulus
generation and analysis so that reaction times could be related to the
momentary position of the approaching sound.

\subsection{Binaural rendering}

Binaural spatialization followed the 3D Tune-In binaural-rendering framework
\citep{CuevasRodriguez2019ThreeDTI}. In simple terms, binaural rendering
transforms a sound source into the two signals that would reach the listener's
left and right ears from a specified position in space. This transformation is
based on head-related impulse responses (HRIRs), which are short acoustic
filters describing how the head, torso, and outer ears shape a sound before it
arrives at each ear. HRIRs provide the spectral and timing differences that
allow a headphone signal to be perceived as coming from a particular direction
and distance rather than from inside the head.

The HRIRs were stored in SOFA format, the Spatially Oriented Format for
Acoustics \citep{Majdak2022SOFA}. SOFA is a standardized way of storing
spatial-acoustics measurements, including the source directions and the
corresponding left- and right-ear impulse responses. The specific HRIR dataset
was FABIAN, a high-resolution measurement of a head-and-torso mannequin
\citep{Brinkmann2017FABIAN}. FABIAN was used as a non-individualized listener
model: the same measured head-and-ear filtering was applied for all
participants. The Python \texttt{sofar} package was used to read and inspect
these SOFA-formatted HRIR data \citep{Sofar2026}.

Rendering was anechoic and distance controlled. Reverberation was disabled, and
direct-path distance attenuation was applied so that the auditory cue changed
systematically as the source approached the listener. The rendered signal was
peak-normalized to preserve digital headroom; this normalization was not an
absolute sound-pressure-level calibration, so acoustic level should be
calibrated at the headphone output before participant testing if absolute SPL
is reported.

\subsection{Spoken breathing instructions}

Each trial began with a spoken instruction specifying the respiratory phase for
that trial. The default protocol used British English spoken instruction
recordings prepared as part of the toolkit stimulus set. The instruction
recordings were normalized to exactly 4.000~s at 44.1~kHz, one channel, and
16-bit resolution. Separate inhale and exhale instructions were used, and each
instruction was matched in duration to the 4.000~s pre-stimulus segment.

\subsection{Tactile stimulation, SOA manipulation, and trial types}

The tactile cue was not an incidental overlay added to selected auditory trials;
the SOA between looming-sound onset and tactile-cue onset was the central
experimental manipulation of the paradigm. Each tactile-present trial therefore
specified both an auditory source and an SOA, so that tactile processing was
sampled at a defined moment after the onset of the task-irrelevant looming
stimulus. The tactile stimulus itself was specified at the software level as a
100~ms vibrotactile pulse. The software tactile waveform is a two-component
sinusoidal pulse with a 200~Hz attack component and a 50~Hz decay component,
shaped by a smooth envelope and normalized to a peak of 0.95 full scale.

For acquisition, the auditory and tactile components were assembled on a
common three-channel audio timeline. The first two channels contained the left
and right binaural auditory signals, and the third channel contained the
vibrotactile signal. On audio--tactile trials, the tactile pulse was placed in
the third channel at the sample corresponding to the scheduled SOA after
looming-sound onset. Thus, the SOA was realized by the temporal structure of
the multichannel stimulus itself rather than by a separate software command
issued during playback. Catch trials retained the same auditory sequence but
contained no tactile pulse.

Within each block, the scheduler generated two balanced trial rows, one for
inhale trials and one for exhale trials. For each respiratory phase and block,
the four auditory spectra were crossed with the five SOAs, yielding 20
audio--tactile trials per phase per block. Three catch trials were added per
phase per block, producing 23 trials for the inhale row and 23 trials for the
exhale row. The two rows were interleaved to form a 46-trial block. Across the
six-block session, this yielded 120 inhale and 120 exhale audio--tactile
trials, plus 18 inhale and 18 exhale catch trials. Trial order was generated
with seed 20250604 using a no-immediate-repeat strategy and a maximum of two
consecutive trials of the same trial type. Participant block order was
counterbalanced by seeded rotation across the 50-participant exported schedule.

\subsection{Session control and event timing}

Before data collection, a participant-specific trial order was generated from
the experimental design. For each trial, the software specified the auditory
source, the respiratory phase, whether a tactile cue was present, and the SOA
for tactile-present trials. The resulting stimulus sequence was presented as a
single synchronized multichannel stream, preserving the sample-level temporal
relationship between the binaural auditory cue and the vibrotactile cue.

Timing-sensitive sessions were run in a native full-screen response interface
rather than in the design interface. Trial events were time-locked to the onset
of audio playback and included trial onset, looming-sound onset, tactile onset
for tactile-present trials, and block boundaries. Behavioural responses were
logged as mouse-click events. An optional Lab Streaming Layer marker stream was
available for synchronized multimodal recordings \citep{Kothe2025LSL}.

\subsection{Hardware-level audio and tactile delivery}

The hardware acquisition route was treated separately from the software
schedule. The preferred route used a single synchronized multichannel audio
interface, with two output channels routed to the headphones and a third output
channel routed to the vibrotactile actuator. This arrangement was chosen to
minimize relative latency between the three stimulus outputs: the two binaural
headphone signals and the tactile signal were emitted by the same audio stream,
through the same hardware device, under the same audio clock. The absolute
audio-interface latency is therefore largely common to all channels, whereas
the experimentally critical quantity is the residual inter-channel skew between
the binaural audio and tactile outputs.

Hardware loopback validation was used as the preferred timing quality-control
procedure for the physical route. The validation routine tested inter-channel
skew, residual jitter, drift, and timing offsets between the headphone and
tactile output channels. This electrical validation establishes whether the
three audio-interface outputs preserve the intended SOAs at the hardware
output. Mechanical actuator latency, however, must be characterized separately
whenever the tactile transducer, wearable actuator, cabling, or amplifier chain
is changed.

\subsection{Reaction-time extraction and quality control}

Immediate quality-control analysis paired each tactile-onset event with the
first behavioural response occurring between 100 and 3000~ms after tactile
onset. Reaction times were summarized by participant, trial type, respiratory
phase, auditory source, and SOA. For task variants including baseline
tactile-only trials, the toolkit can compute SOA-matched baseline-corrected
reaction times. For the present Study 5 preload, which omitted baseline trials,
the primary exported quality-control metric was the mean reaction time at each
SOA and condition. When at least four SOA levels were available, the software
also fitted a sigmoid function to the reaction-time profile using SciPy
\citep{Virtanen2020SciPy}, following the convention that the fitted transition
point indexes the spatial boundary at which task-irrelevant auditory
information begins to facilitate tactile processing
\citep{Canzoneri2012DynamicSounds,Serino2015BodyPart}. Final inferential
models, participant exclusions, and any preregistered analysis choices should
be reported in the statistical-analysis section of the manuscript.
